\chapter{Research Chronology}
\label{app:chronology}
\section{Original Objective and Dataset Selection}
The original research objective concerned process-aware multimodal architectures for
robotic welding. The exposé defined the research questions before the subsequent
industrial and robotic studies. The choice and timing of those studies were exploratory.
Their relationship to the original questions is a retrospective
interpretation.\ev{Original scope and scientific chronology}{\EvidenceRef{E002}{pp.
2--3}, \EvidenceRef{E004}{sections 10--11}}

\section{From the Welding Motivation to Two Task-Specific Studies}
The PHM study began with image exploration and developed a continuous gear-condition
proxy. The September v3 target used one deterministic view per tooth, teeth 5--28 and a
top-three run aggregate. The earlier maximum-over-views targets and fixed image split
addressed different prediction settings. The available records do not establish the
initial reason for selecting PHM. Concerns about the validity of the proxy emerged
through its construction and evaluation.\ev{PHM target
chronology}{\EvidenceRef{E004}{sections 12--14}, \EvidenceRef{E023}{}}

Re-examination of Intel Welding found intended-condition categories rather than
independently verified defects, class-dependent modality availability, recipe/category
confounding and 216 sessions crossing the supplied split. A proposed 32-session holdout
had already appeared in a whole-population context audit. That exposure invalidates an
untouched-test claim without proving that every model trained on those labels. The
foundation review stopped before new model training; it supplies feasibility evidence,
not a negative fusion result. PHM was already under way. REASSEMBLE subsequently
provided native completed-action outcomes, aligned observations and recording identities
for comparative evaluation; unresolved waveform timing and spoken outcome leakage
prevented an acoustic predictor.\ev{Welding and assembly
eligibility}{\EvidenceRef{E120}{}, \EvidenceRef{E004}{sections 14--15}}

The September PADERBORN validation is retained as a bounded investigation, not an
additional principal case study. Severity was fixed per bearing. The requested
condition-only comparison favored the vibration predictor over the constant, whereas the
stricter bearing-and-condition contrast was inconclusive: MAE difference $-.1184$ [95\%
interval $-.2978$, .0589]. The report promoted the stricter protocol as primary without
resolving that protocol fork. Recoverable bearing identity makes a shortcut plausible,
not proven as the sole mechanism. Gate~3 distillation was not run; withholding current
as privileged information was a study choice, not absence of current from the
dataset.\ev{PADERBORN protocol fork and stopping boundary}{\EvidenceRef{E117}{}}

\section{Follow-Ups, Human Review and External Cases}
September follow-ups examined longer sensor training, a same-sensor hybrid, anchored
adaptive coefficients, matched augmentation and boundary-aware visual pooling. These
include positive exploratory results, but they did not execute a combined system
containing every favorable choice. The discontinued temporal-model work is distinct from
completed pooling experiments.

The October experiments examined a spatial hybrid encoder, fine-tuning with unequal
training budgets, a matched two-epoch comparison and noise/quality interactions on
selected folds. GPU availability limited the 100-epoch comparison to one previously
inspected assessment fold. The later F2 fit also used a different set of training
predictions from the original experiment. The results therefore describe a follow-up on
that fold.\ev{Follow-up chronology and scope}{\EvidenceRef{E004}{sections 18--21},
\EvidenceRef{E003}{sections 3.3 and 4.4}}

A reviewer used the September interface to assess 64 REASSEMBLE clips. The returned
ratings identified disagreements and potential boundary or presentation problems.
Adjudication was incomplete, and the evaluation used the original labels.

The earlier \enquote{Milling Dataset} workspace and MSM refer to the same source, not
independent datasets. The 114-case feasibility pilot failed positive-group support;
later eligibility and alignment work admitted 1,895 intervals and enabled comparisons
under several grouping and label-2 policies. Thus, a failed pilot gate did not establish
failure of the later study. HANDOVER's pre-score common-window amendment retained all
151 official test trials, but absent participant mapping prevented the mandated OOF bank
for learned fusion. These stopping and continuation decisions concern different
requirements.\ev{Milling identity and external eligibility}{\EvidenceRef{E121}{},
\EvidenceRef{E108}{}}

REBOOT remained a feasibility investigation. The 28 October decision admitted no
normal/abnormal windows or supervised groups under the requested protocol despite 60
annotated cylinder-install episodes. Visual review was pending; no CNN extraction,
training or performance estimate followed. The wider screening excluded RoboFAC for lack
of real numerical counterparts, retained FAILURE only as a proposed RGB--audio pilot
addressing a different input question, and removed CRIA at the researcher's request.
Screening outcomes are not model results.\ev{REBOOT and screened
alternatives}{\EvidenceRef{E118}{}, \EvidenceRef{E119}{}}

ViFailback progressed from cohort and telemetry checks to original frozen evaluation,
matched frozen/donor representation studies and bounded target-CNN adaptation. The 24
November closure records 79 of 111 planned fits complete. FULL20 has one seed and one
outer fold with all three required inner fits; FULL100 is a separate pilot. Their paired
adaptation effects remain inconclusive. The researcher stopped the remaining 32 FULL20
fits, and corruption evaluation remained unexecuted. This revision extends reporting
through that closure while preserving the earlier protocols and complete frozen/donor
results.\ev{ViFailback extended cutoff}{\EvidenceRef{E115}{}, \EvidenceRef{E116}{}}

